\begin{exercise}{55}
Show that every simplicial set $A$ with only finitely many non-degenerate
simplices is \emph{compact}, i.e., the canonical map
\[
 \operatorname*{colim}_{i\in I}\mathrm{Hom}_{\mathrm{sSet}}(A,X_i)
 \longrightarrow
 \mathrm{Hom}_{\mathrm{sSet}}\Bigl(A,\operatorname*{colim}_{i\in I}X_i\Bigr)
\]
is an isomorphism, provided $I$ is a filtered category.

\begin{proof}
We first recall how equality works in filtered colimits of sets.
In the quotient construction of Exercise~\ref{ex:6}, two
representatives are equal precisely when they have equal images
at a common stage. Common-stage equality implies equality
in the quotient. Conversely, the common-stage relation is reflexive,
symmetric, and contains the generating identifications.
For transitivity, send two witness stages to a common target
and equalize the two induced arrows from the shared middle stage,
using Exercise~\ref{ex:53}.
If both representatives already lie at the same stage,
equalizing the two witness arrows shows that a single transition
makes them equal. Iterating this observation makes any finite
collection of such equalities hold after a single transition.

Put $X=\operatorname*{colim}_iX_i$, computed degreewise by
Proposition~1.1.23, and let $D$ be the finite set of nondegenerate
simplices of $A$.
By Exercises~\ref{ex:3} and~\ref{ex:5}, assigning images $z_\sigma$
to $\sigma\in D$ defines a simplicial map exactly when
\[
 \begin{gathered}
  \iota^*z_\sigma=\beta^*z_\tau
  \quad\text{whenever}\quad \iota^*\sigma=\beta^*\tau,\\
  \iota\text{ injective},\quad
  \beta\text{ surjective},\quad\tau\in D.
 \end{gathered}
\]
There are only finitely many such conditions: for each $\sigma$,
there are finitely many injections into its indexing ordinal,
and Exercise~\ref{ex:5} determines $(\beta,\tau)$ uniquely.
Indeed, extend the assignment by
$z_{\alpha^*\sigma}=\alpha^*z_\sigma$ for surjective $\alpha$.
To check a further pullback, factor its composite with $\alpha$
as a surjection followed by an injection and apply the displayed
condition. This proves compatibility with every simplicial operator;
uniqueness follows from Exercise~\ref{ex:5}.

Given $f\colon A\to X$, choose representatives for the finitely
many values $f(\sigma)$, and send their stages to one stage $i$
using Exercise~\ref{ex:53}.
The displayed conditions hold in $X$.
The equality criterion above, applied degreewise and finitely
many times, makes them all hold after one transition $i\to j$.
The resulting assignment defines $A\to X_j$ inducing $f$.
This proves surjectivity.

For injectivity, suppose $A\to X_i$ and $A\to X_j$ induce
the same map to $X$.
Send $i,j$ to a common stage.
Their values on the finite set $D$ become equal after a
further single transition by the same equality criterion.
The two maps then agree on every simplex by Exercise~\ref{ex:5},
so represent the same element of the colimit of hom-sets.
If $A$ is empty, the same argument uses the nonemptiness
of $I$, which is included in Exercise~\ref{ex:53}.
\end{proof}
\end{exercise}
